\section{Introduction}

A bijective sliding block code between subshifts has a sliding-block inverse,
but compactness alone does not make the inverse radius uniform over a class of
systems.  This paper isolates that loss of uniformity for perhaps the simplest
nontrivial local map on a binary alphabet,
\begin{equation}\label{eq:xor-intro}
  \Delta(x)_i=x_i+x_{i+1}\pmod 2.
\end{equation}
The forward map has memory $0$ and anticipation $1$.  We construct a family on
which it is one-to-one while the least symmetric radius of its inverse tends
to infinity.

Rank-one symbolic systems provide a controlled setting for this question.
Their word-building language and topological isomorphisms were developed in
particular by Gao and Hill~\cite{GaoHill2016}; Gao and Ziegler proved that a
subshift factor of a rank-one subshift is either finite or isomorphic to the
source~\cite{GaoZiegler2020}.  Finite spacer rank and its language-theoretic
scope are treated by Gao, Jacoby, Johnson, Leng, Li, Silva, and
Wu~\cite{GaoEtAl2025}.  Those qualitative results motivate, but are not used
to replace, the quantitative inverse-window proof below.  In another
structured class, Durand and Leroy obtain computable radius bounds for maps
between two prescribed minimal substitution subshifts~\cite{DurandLeroy2022}.
Our question is different: whether coarse uniform rank-one bounds control an
inverse radius.  The answer for the explicit XOR code is negative.

We first record a finite-group principle.  For a subshift $X\subseteq G^\Z$,
let $cX$ denote coordinatewise left multiplication by $c\in G$, and write
$\Lang_\ell(X)$ for its words of length $\ell$.  Define
\begin{equation}\label{eq:separation-length-intro}
 s_G(X)=\min\bigl\{\ell\geq 1:
 \Lang_\ell(X)\cap c\Lang_\ell(X)=\varnothing
 \text{ for every }c\neq e\bigr\},
\end{equation}
with $s_G(X)=\infty$ if the set is empty.  The exact result is as follows.

\begin{theorem}[finite-group derivative criterion]\label{thm:finite-group-intro}
Let $G$ be finite, $X\subseteq G^\Z$ a nonempty subshift, and
$\Delta_G(x)_i=x_i^{-1}x_{i+1}$.  Then $\Delta_G|_X$ is injective if and only
if $X\cap cX=\varnothing$ for every $c\neq e$.  In that case $s_G(X)<\infty$,
$\Delta_G:X\to \Delta_G(X)$ is a conjugacy, and its inverse can use the window
$[-m,n]$, where $m,n\geq0$, if and only if
\begin{equation}\label{eq:window-intro}
 m+n+2\geq s_G(X).
\end{equation}
In particular its least symmetric inverse radius is
\begin{equation}\label{eq:radius-intro}
 R_{\rm inv}(X)=\min\{r\geq0:2r+2\geq s_G(X)\}
 =\left\lceil\frac{s_G(X)-2}{2}\right\rceil.
\end{equation}
\end{theorem}

The family giving unbounded radii has a short definition.  For each $N\geq1$
put $v^{(N)}_0=0$ and
\begin{equation}\label{eq:construction-intro}
 v^{(N)}_{n+1}=
 \begin{cases}
  v^{(N)}_n1v^{(N)}_n,&0\leq n<N,\\
  v^{(N)}_nv^{(N)}_nv^{(N)}_n11v^{(N)}_n,&n\geq N.
 \end{cases}
\end{equation}
The words are nested prefixes.  Let $V_N$ be their one-sided limit and $X_N$
the two-sided subshift whose language consists of the factors of $V_N$.

\begin{theorem}[bounded rank-one family]\label{thm:main-intro}
For every $N\geq1$, the subshift $X_N$ is nonempty, minimal, aperiodic, and
uniquely ergodic.  Its unique measure has one-density
\begin{equation}\label{eq:density-intro}
 \rho_N=\frac{3\cdot2^N-1}{3\cdot2^{N+1}-1}<\frac12.
\end{equation}
The XOR derivative in \eqref{eq:xor-intro} is therefore a conjugacy
$X_N\to\Delta(X_N)$.  If $R_N$ is the least symmetric radius of its inverse,
then
\begin{equation}\label{eq:main-bound-intro}
 s_{\Ftwo}(X_N)=3|v_N|+2=3\cdot2^{N+1}-1,
 \qquad R_N=3\cdot2^N-1.
\end{equation}
Across the family, the alphabet has size $2$, cutting numbers are at most
$4$, spacer lengths are at most $2$, and the forward code is fixed.
\end{theorem}

The proof has four independent pieces.  The finite-group fibre calculation
gives the exact inverse-window criterion.  Bounded-gap copies of each
construction word give minimality, while arbitrarily long right-special
words give aperiodicity.  A direct occurrence-count estimate proves uniform
word frequencies and yields the exact density.  Finally, a word of length
$3|v_N|+1$ and its bitwise complement both occur in $V_N$.  Coding equal
adjacent pairs as typed defects shows that a longer complementary pair is
impossible: the tail spacer language avoids the pattern $101$.  This proves
the exact separation and radius formulas in \eqref{eq:main-bound-intro}.
